\documentclass[11pt]{article}
\usepackage[margin=1in]{geometry}
\usepackage{amsmath,amssymb,amsthm}
\usepackage{enumitem}
\usepackage{xcolor}
\usepackage{microtype}
\usepackage[hidelinks]{hyperref}
\emergencystretch=2em

\newtheorem{theorem}{Theorem}
\newtheorem{lemma}{Lemma}
\newtheorem{proposition}{Proposition}
\theoremstyle{definition}
\newtheorem*{assumptionA}{Assumption A}
\newtheorem*{assumptionB}{Assumption B}
\theoremstyle{remark}
\newtheorem{remark}{Remark}

\newcommand{\R}{\mathbb{R}}
\newcommand{\E}{\mathbb{E}}
\newcommand{\pr}{\mathrm{pr}}
\newcommand{\TV}{\mathrm{TV}}
\newcommand{\Dr}{D^{(r)}}
\newcommand{\lamtv}{\lambda_{\TV}}
\newcommand{\thetas}{\theta^{*}}
\newcommand{\en}{\varepsilon_n}
\newcommand{\Lh}{\widehat{\mathcal L}}
\newcommand{\Lm}{\mathcal L}
\newcommand{\Rh}{\widehat{\mathcal R}}
\newcommand{\Ql}{Q_\lambda}
\newcommand{\Pp}{P^{\perp}}
\newcommand{\cP}{\mathcal P}
\newcommand{\fl}{\underline f}
\newcommand{\fu}{\bar f}
\newcommand{\Vb}{\bar V}
\newcommand{\nn}[1]{\|#1\|_n}
\newcommand{\ind}{\mathbf 1}
\newcommand{\op}{\mathrm{op}}

\title{Posterior contraction for Bayesian quantile trend filtering}
\author{}
\date{24 November 2026}

\begin{document}
\maketitle
\vspace{-2em}

%======================================================================
\section{Model}

Fix a quantile level $\tau\in(0,1)$ and an integer $r\ge2$ (trend filtering of order $k=r-1$; the default $k=2$ is $r=3$). We observe independent $y_1,\dots,y_n$ with distribution functions $F_{y_i}$, and $\thetas_i$ denotes the $\tau$-quantile of $y_i$. With $D^{(1)}\theta=(\theta_2-\theta_1,\dots,\theta_n-\theta_{n-1})$ and $\Dr=D^{(1)}D^{(r-1)}\in\R^{m\times n}$, $m=n-r$, let
\[
\TV^{(r)}(\theta)=n^{r-1}\|\Dr\theta\|_1,\qquad V^{*}=\TV^{(r)}(\thetas),\qquad
\Delta_n^2(v)=\frac1n\sum_{i=1}^n\min\{|v_i|,v_i^2\},
\]
and $\rho_\tau(u)=u\{\tau-\ind(u<0)\}$. The working likelihood is the asymmetric Laplace density $p(y_i\mid\theta_i)=\tau(1-\tau)\sigma^{-1}\exp\{-\rho_\tau(y_i-\theta_i)/\sigma\}$ and the prior is $p(\theta\mid\lambda)\propto\lambda^m\exp\{-\lambda\|\Dr\theta\|_1\}$, which is flat on polynomials of degree at most $r-1$. We assume $\sigma=1$ for simplicity; the theorem holds for any fixed $\sigma=\sigma_0>0$ (end of Theorem~\ref{thm:main}). The tuning $\lambda$ is deterministic. The posterior is, for Borel $B\subset\R^n$,
\begin{equation}\label{eq:post}
\Pi_\lambda(B\mid y)=\frac{\int_B\exp\{-\sum_{i}\rho_\tau(y_i-\theta_i)-\lambda\|\Dr\theta\|_1\}\,d\theta}
{\int_{\R^n}\exp\{-\sum_{i}\rho_\tau(y_i-\theta_i)-\lambda\|\Dr\theta\|_1\}\,d\theta}.
\end{equation}
It is proper (Lemma~\ref{lem:coord}). Equivalently the penalty is $\lamtv\TV^{(r)}(\theta)$ with $\lamtv=\lambda/n^{r-1}$.

%======================================================================
\section{Assumptions}

\begin{assumptionA}[Madrid Padilla and Chatterjee, 2022]
There exist constants $L>0$ and $\fl>0$ such that for any $n$, any $\delta\in\R^n$ with $\|\delta\|_\infty\le L$ and all $i$,
$\bigl|F_{y_i}(\thetas_i+\delta_i)-F_{y_i}(\thetas_i)\bigr|\ge\fl\,|\delta_i|$.
\end{assumptionA}
\begin{assumptionB}
There exists a constant $\fu<\infty$ such that, for the same $L$, any $n$, any $\|\delta\|_\infty\le L$ and all $i$,
$\bigl|F_{y_i}(\thetas_i+\delta_i)-F_{y_i}(\thetas_i)\bigr|\le\fu\,|\delta_i|$.
\end{assumptionB}

\noindent Assumption~B excludes an atom at the quantile, so $F_{y_i}(\thetas_i)=\tau$. No moments are assumed (Cauchy noise is allowed). Finally, $V^{*}\le\Vb$ for all $n$.

%======================================================================
\section{Theorem}

\begin{theorem}\label{thm:main}
Let $y_1,\dots,y_n$ be independent and satisfy Assumptions A and B, and let $V^{*}\le\Vb$ for all $n$. Let
\[
\lambda=c\,n^{(2r-1)/2},\qquad\text{equivalently}\qquad\lamtv=c\,n^{1/2},
\]
with a constant $c>0$. There exist constants $M$, $C$, $c_0>0$ and $n_0$, depending only on $c$, $r$, $L$, $\fl$, $\fu$ and $\Vb$, such that for all $n\ge n_0$
\begin{align*}
&\text{(a)}\quad \E_{\thetas}\,\Pi_\lambda\bigl(\theta:\ \Delta_n^2(\theta-\thetas)>M\,n^{-(2r-1)/(2r)}\bigm|y\bigr)\le C\,n^{-1/(2r)},\\
&\text{(b)}\quad \E_{\thetas}\,\Pi_\lambda\bigl(\theta:\ \Delta_n^2(\theta-\thetas)\le c_0\,n^{-(2r-1)/(2r)}\bigm|y\bigr)\le 2\,n^{-1/(2r)}.
\end{align*}
The bounds are uniform over all laws of $y$ that satisfy A and B with the given constants and $V^{*}\le\Vb$. The same holds for any fixed $\sigma=\sigma_0>0$, with constants depending also on $\sigma_0$.
\end{theorem}

\noindent So the posterior contracts at the rate $n^{-(2r-1)/(2r)}$, and by (b) this rate is exact: no logarithm, no smaller power. For the default $r=3$: $\lambda=c\,n^{5/2}$ and the rate is $n^{-5/6}$.

\begin{remark}[Why not $n^{-2r/(2r+1)}$]
Madrid Padilla and Chatterjee (2022) obtain $n^{-2r/(2r+1)}$, up to a logarithm, for the estimator. The whole posterior is slower: it averages over windows of $w\asymp\lambda^{1/r}$ points, so its own spread is $\Delta_n^2\asymp\lambda^{-1/r}$, while its smoothing bias is of order $V^{*2}\lambda^{(2r-1)/r}n^{-(2r-1)}$. The two balance at $\lambda\asymp n^{(2r-1)/2}$, giving $n^{-(2r-1)/(2r)}$, the exponent of the lower bound for linear estimators over this class (Donoho and Johnstone, 1998; Tibshirani, 2014, Section~5.1 with $k=r-1$).
\end{remark}

\begin{remark}[Misspecification]
The asymmetric Laplace density is a working likelihood: we do not assume that the data follow it, and every probability statement in Theorem~\ref{thm:main} is under the true distribution of the data. The posterior is therefore a Gibbs, or quasi-, posterior built on the check loss, in which $1/\sigma$ acts as a learning rate (Syring and Martin, 2023; Atchad\'e, 2017). It concentrates on the true $\tau$-quantile because, whatever the distribution of $y_i$, the $\tau$-quantile minimizes the expected check loss and hence the Kullback--Leibler divergence from the truth to the asymmetric Laplace family (Sriram, Ramamoorthi and Ghosh, 2013). Misspecification changes the spread of the posterior and not its centre, so the theorem describes where the posterior concentrates and how fast, not the frequentist coverage of its credible sets (Yang, Wang and He, 2016).
\end{remark}

\begin{remark}[Learned $\sigma$]
In the sampler $\sigma\sim\mathrm{IG}(a,b)$, so $\sigma\mid\theta,y\sim\mathrm{IG}(n+a,\,S(\theta)+b)$ with $S(\theta)=\sum_i\rho_\tau(y_i-\theta_i)$. Hence $\sigma$ concentrates near $S(\theta)/n$, close to the average check loss $\bar\sigma_n=n^{-1}\sum_i\E\rho_\tau(y_i-\thetas_i)$ when this is finite and bounded away from zero (not for Cauchy noise). The theorem with $\sigma_0=\bar\sigma_n$ describes the posterior at that scale; a learned-$\sigma$ version is not claimed.
\end{remark}

\noindent\emph{The mode.} The posterior mode is quantile trend filtering with tuning $\sigma_0\lamtv$; at $\lamtv=c\,(n\log n)^{1/(2r+1)}$ every mode attains $\Delta_n^2\le C\,n^{-2r/(2r+1)}(\log n)^{1/(2r+1)}$ with probability at least $1-C(n\log n)^{-1/(2r+1)}$, by the localization argument of Proposition~\ref{prop:local} (proof not given here).

%======================================================================
\section{Outline of the proof}\label{sec:outline}

Write $\delta=\theta-\thetas$, $Z_i=y_i-\thetas_i$, $\nn{v}=n^{-1/2}\|v\|_2$, and
\[
\Lh(\delta)=\sum_{i=1}^n\{\rho_\tau(Z_i-\delta_i)-\rho_\tau(Z_i)\},\qquad \Lm(\delta)=\E_{\thetas}\Lh(\delta),\qquad \en^2=n^{-(2r-1)/(2r)}.
\]
Then $n\en^2=n^{1/(2r)}$ and $(n\en^2)^{2r}=n$. The posterior is a ratio $\Pi_\lambda(B\mid y)=N(B)/N(\R^n)$, $N(B)=\int_Be^{-\Lh(\theta-\thetas)}d\nu(\theta)$, with $\nu$ the prior. Five steps:
\begin{enumerate}[itemsep=2pt,topsep=3pt,label=\textbf{\arabic*.},leftmargin=*]
\item \emph{Coordinates.} $\theta=Nc+Su$ with $u=\Dr\theta$: under the prior $c\in\R^r$ is flat and the $u_j$ are independent Laplace$(\lambda)$.
\item \emph{Denominator.} With high probability $N(\R^n)\ge(\text{volume})\cdot e^{-Cn\en^2}$. This needs prior mass $e^{-Cn\en^2}$ near the truth. Centring at $\thetas$ costs $\lambda\|\Dr\thetas\|_1=cV^{*}\sqrt n$, in general far above $n\en^2$, so the ball is centred at a locally smoothed truth, at a quadratic cost $\lambda^2\|u^\circ\|_2^2\lesssim n\en^2$; the centred small-ball probability is bounded without a logarithm.
\item \emph{Localization.} The posterior puts exponentially small mass outside a set $K_n$ on which $\|\delta\|_\infty\to0$. This is where $r\ge2$ is needed. On $K_n$, $\Delta_n^2(\delta)=\nn{\delta}^2$ and Assumptions A and B apply to every coordinate.
\item \emph{Shells.} On a shell $\{a\en^2<\nn{\delta}^2\le2a\en^2\}$ write $\Lh(\delta)=-\langle\psi,\delta\rangle+\Rh(\delta)$ with the score $\psi_i=\tau-\ind(Z_i<0)$ and a nonnegative remainder. A spline projector with $J\asymp Mn\en^2$ knots splits the score: the coarse part is handled by Cauchy--Schwarz, the fine part by the exact Laplace moment generating function. The numerator is of order $e^{-(\fl/8)an\en^2}$.
\item \emph{Assembly.} $\E\,\Pi\le\pr(\text{bad events})+\E[\Pi\,\ind_{\rm good}]$; the volume factors cancel. The bad events have probability $O(1/(n\en^2))=O(n^{-1/(2r)})$.
\end{enumerate}
\emph{Where the rate comes from.} The prior mass of an $\en$-ball is $\exp\{-C(c\en)^{-2/(2r-1)}\}$, which must be at least $e^{-Cn\en^2}$; and the fine score costs $n/(c^2(Mn\en^2)^{2r-1})$ in the exponent, which must be at most of order $Mn\en^2$. Both hold, up to constants, exactly at $\en^2=n^{-(2r-1)/(2r)}$. For (b), the prior mass of a ball of radius $\sqrt{c_0}\,\en$ is too small for small $c_0$.

%======================================================================
\section{Proof}

\paragraph{Reduction to $\sigma=1$.} For $\sigma_0>0$, the substitution $\theta=\sigma_0\theta'$ and positive homogeneity of $\rho_\tau$ show that the law of $\theta/\sigma_0$ under the $\sigma_0$-posterior is the $\sigma=1$ posterior based on $y/\sigma_0$ with $c$ replaced by $\sigma_0c$. The data $y/\sigma_0$ satisfy A and B with $(L/\sigma_0,\sigma_0\fl,\sigma_0\fu)$, and $V^{*}\le\Vb/\sigma_0$. Since $\min\{\sigma_0,\sigma_0^2\}\Delta_n^2(v)\le\Delta_n^2(\sigma_0v)\le\max\{\sigma_0,\sigma_0^2\}\Delta_n^2(v)$, (a) and (b) for $\sigma_0$ follow from the case $\sigma=1$ applied to the rescaled problem, with its $M$ multiplied by $\max\{\sigma_0,\sigma_0^2\}$ and its $c_0$ by $\min\{\sigma_0,\sigma_0^2\}$. From now on $\sigma=1$ and $\lambda=c\,n^{r-1/2}$.

%----------------------------------------------------------------------
\subsection{Coordinates and the check loss}

For integers $t$ and $p\ge0$ let $b_p(t)=\binom tp\ind(t\ge0)$, with $\binom tp=t(t-1)\cdots(t-p+1)/p!$. Pascal's rule gives
\begin{equation}\label{eq:pascal}
b_p(t+1)-b_p(t)=b_{p-1}(t)\quad(p\ge1),\qquad b_0(t+1)-b_0(t)=\ind(t=-1).
\end{equation}
Let $S\in\R^{n\times m}$ have columns $s_j(i)=b_{r-1}(i-j-1)$ (so $s_j(i)=0$ for $i<j+r$ and $0\le s_j(i)\le n^{r-1}/(r-1)!$). Let $\cP\subset\R^n$ be the vectors $(p(i))_{i\le n}$ with $p$ a polynomial of degree at most $r-1$, $N\in\R^{n\times r}$ an orthonormal basis of $\cP$, $P=NN^\top$, $\Pp=I-P$. Let $\Ql$ be the law on $\R^m$ of $m$ independent Laplace variables with density $(\lambda/2)e^{-\lambda|x|}$, and $V_r$ the volume of the unit ball of $\R^r$.

\begin{lemma}[coordinates]\label{lem:coord}
Let $n>r$ and $D=\Dr$.
\begin{enumerate}[label=(\alph*),itemsep=1pt,topsep=2pt]
\item $DS=I_m$ and $\ker D=\cP$.
\item $(c,u)\mapsto Nc+Su$ is a linear bijection $\R^r\times\R^m\to\R^n$, with $u=D\theta$. Let $\nu$ be the image of $\mathrm{Leb}_r\otimes\Ql$. Then
\begin{equation}\label{eq:ratio}
\Pi_\lambda(B\mid y)=\frac{N(B)}{N(\R^n)},\qquad N(B)=\int_Be^{-\Lh(\theta-\thetas)}\,d\nu(\theta),\qquad 0<N(\R^n)<\infty .
\end{equation}
\item Let $u^{*}=D\thetas$ and $c^{*}=N^\top(\thetas-Su^{*})$. For $\theta=Nc+Su$ and $\delta=\theta-\thetas$,
\begin{equation}\label{eq:pyth}
\|\delta\|_2^2=\|c-\bar c(u)\|_2^2+\|\Pp S(u-u^{*})\|_2^2,\qquad \bar c(u)=c^{*}-N^\top S(u-u^{*}),
\end{equation}
and for measurable $H\ge0$ and $\varrho\ge0$ on $\R^m$,
\begin{equation}\label{eq:vol}
\int\ind\{\|c-\bar c(u)\|_2\le\varrho(u)\}H(u)\,d\nu=V_r\int\varrho(u)^rH(u)\,d\Ql(u);
\end{equation}
in particular $\nu\{\|\delta\|_2\le\varrho\}\le V_r\varrho^r$ for every constant $\varrho\ge0$.
\end{enumerate}
\end{lemma}

\begin{proof}
(a) $D$ is the $r$-th forward difference. Applying \eqref{eq:pascal} $r$ times to $i\mapsto b_{r-1}(i-j-1)$ gives $\ind(i=j)$, so $DS=I_m$ and $\operatorname{rank}D=m$. $D$ annihilates $\cP$ and $\dim\cP=r=n-m$, so $\ker D=\cP$.
(b) If $Nc+Su=0$, applying $D$ gives $u=0$, then $c=0$; so the map is bijective, with constant Jacobian, and $e^{-\lambda\|u\|_1}du=(2/\lambda)^md\Ql(u)$. Substitute in \eqref{eq:post} and divide both integrals by $\exp\{-\sum_i\rho_\tau(Z_i)\}$. For finiteness: $\rho_\tau(x)\ge\min\{\tau,1-\tau\}|x|$ and $\|y-Nc-Su\|_1\ge\|c-N^\top(y-Su)\|_2$, so the $c$-integral is finite and free of $u$.
(c) $\thetas-Su^{*}\in\ker D=\cP$, so $\thetas=Nc^{*}+Su^{*}$, $P\delta=N\{c-\bar c(u)\}$ and $\Pp\delta=\Pp S(u-u^{*})$; \eqref{eq:pyth} is Pythagoras and \eqref{eq:vol} is Fubini.
\end{proof}

Let $F_i(t)=\pr(Z_i\le t)$, $\ell_i(t)=\rho_\tau(Z_i-t)-\rho_\tau(Z_i)$, $m_i(t)=\E\ell_i(t)$, $\psi_i=\tau-\ind(Z_i<0)$, so $\Lh(\delta)=\sum_i\ell_i(\delta_i)$, $\Lm(\delta)=\sum_im_i(\delta_i)$. Put $C_1=\max\{\fu/2,1/L\}$.

\begin{lemma}[check loss]\label{lem:check}
Under Assumptions A and B, for every $i$ and all real $t,t'$:
\begin{enumerate}[label=(\alph*),itemsep=1pt,topsep=2pt]
\item $\ell_i$ is convex, $\ell_i(0)=0$, $|\ell_i(t)-\ell_i(t')|\le|t-t'|$, and $\ell_i(t)$ takes values in a fixed interval of length $|t|$.
\item $F_i(0)=\tau$; the $\psi_i$ are i.i.d., with values in $\{\tau-1,\tau\}$, mean $0$ and variance $\tau(1-\tau)\le1/4$.
\item (Knight's identity) $\ell_i(t)=-t\psi_i+R_i(t)$, where $R_i(t)=(t-Z_i)\ind(0\le Z_i<t)$ for $t>0$, $R_i(t)=(Z_i-t)\ind(t<Z_i<0)$ for $t<0$, $R_i(0)=0$; so $0\le R_i(t)\le|t|$.
\item $m_i(t)=\E R_i(t)=\int_0^t\{F_i(s)-\tau\}ds\ge0$, with $\tfrac12\fl t^2\le m_i(t)\le\tfrac12\fu t^2$ for $|t|\le L$ and $m_i(t)\le C_1t^2$ for all $t$.
\item $\E R_i(t)^2\le\fu|t|^3/3$ for $|t|\le L$, and $\E e^{-2R_i(t)}\le e^{-m_i(t)}$ for $|t|\le\min\{L,3\fl/(4\fu)\}$.
\item $\E e^{-\ell_i(t)}\le e^{t^2/8}$ for all $t$.
\end{enumerate}
\end{lemma}

\begin{proof}
(a), (c). $\rho_\tau$ is convex and $1$-Lipschitz. For $t>0$ the cases $Z_i<0$, $0\le Z_i<t$, $Z_i\ge t$ give $\ell_i(t)=(1-\tau)t$, $(1-\tau)t-Z_i$, $-\tau t$ and $\psi_i=\tau-1,\tau,\tau$, so $\ell_i(t)+t\psi_i=0,\ t-Z_i,\ 0$, and $\ell_i(t)\in[-\tau t,(1-\tau)t]$; $t<0$ is symmetric (Knight, 1998).
(b) By B, $F_i$ is continuous at $0$; since $\thetas_i$ is a $\tau$-quantile, $F_i(0)\ge\tau\ge F_i(0-)$, so $F_i(0)=\tau$ and $\pr(\psi_i=\tau-1)=\tau$ for all $i$.
(d) $\E\psi_i=0$, so $m_i=\E R_i$; for $t>0$, $R_i(t)=\int_0^t\ind(0\le Z_i<s)ds$ gives the integral (Belloni and Chernozhukov, 2011), and similarly for $t<0$. Under A and B, $\fl|s|\le|F_i(s)-\tau|\le\fu|s|$ for $|s|\le L$, with the sign of $s$. For $|t|>L$, (a) gives $m_i(t)\le|t|\le t^2/L$.
(e) $R_i(t)^2=\int_0^t2(t-s)\ind(0\le Z_i<s)ds$ gives $\E R_i^2\le\fu\int_0^t2(t-s)s\,ds=\fu t^3/3$. With $e^{-x}\le1-x+x^2/2$ ($x\ge0$),
$\E e^{-2R_i}\le1-2m_i+\tfrac23\fu|t|^3\le1-m_i\le e^{-m_i}$, using $\tfrac23\fu|t|^3\le\tfrac12\fl t^2\le m_i$.
(f) Hoeffding's lemma with (a) and $m_i\ge0$.
\end{proof}

%----------------------------------------------------------------------
\subsection{Spline approximation}

For an integer $1\le J\le m$ let $h=\lceil m/J\rceil$, $J'=\lfloor m/h\rfloor$, $V_J=\cP+\mathrm{span}\{s_{lh}:l=1,\dots,J'\}$ (discrete splines with $J'$ equally spaced knots), and $P_J$ the orthogonal projector onto $V_J$. Then $h\le2n/J$ and $\dim V_J\le J+r$.

\begin{lemma}[spline approximation]\label{lem:F}
Let $1\le J\le m$. For every $j$ there is $e_j$ with $s_j-e_j\in V_J$, supported on at most $rh$ consecutive indices in $[j-rh,j+rh]$, with $\|e_j\|_\infty\le c_\infty h^{r-1}$, $c_\infty=(1+2^{r-1})(2r)^{r-1}/(r-1)!$. Consequently $(I-P_J)N=0$ and, with $c_F=\sqrt r\,c_\infty$, $c_{\op}=\sqrt3\,r\,c_\infty$,
\begin{equation}\label{eq:F}
\max_{j}\|(I-P_J)s_j\|_2\le c_F\,h^{r-1/2},\qquad \|(I-P_J)S\|_{\op}\le c_{\op}\,h^{r}.
\end{equation}
\end{lemma}

\begin{proof}
For fixed $i$, $p_i(x)=\binom{i-x-1}{r-1}$ is a polynomial of degree $r-1$ with $s_l(i)=p_i(l)$ for $l\le i-1$ and $s_l(i)=0$ for $l\ge i$. For fixed $j$, $q_j=(\binom{i-j-1}{r-1})_i\in\cP$, $s_j(i)=q_j(i)\ind(i\ge j+1)$ and $|q_j(i)|\le(j-i+r-1)^{r-1}/(r-1)!$ for $i\le j$.

\emph{Case $J'\ge r$, $h\le j\le J'h$.} Let $a=\min\{\lfloor j/h\rfloor,J'-r+1\}$, knots $\kappa_l=(a+l)h$, $l=0,\dots,r-1$, $x=j/h-a\in[0,r-1]$, and $a_l$ the Lagrange weights for the nodes $0,\dots,r-1$ at $x$. Put $e_j=s_j-\sum_la_ls_{\kappa_l}$. For $i\le\kappa_0$ all terms vanish; for $i>\kappa_{r-1}$ all are values of $p_i$ and $e_j(i)=p_i(j)-\sum_la_lp_i(\kappa_l)=0$. So $e_j$ lives on $(r-1)h$ consecutive indices, where all entries are at most $((r-1)h)^{r-1}/(r-1)!$. Since $\prod_{k\ne l}|x-k|\le(r-1)!$, $|a_l|\le\binom{r-1}{l}$ and $\sum_l|a_l|\le2^{r-1}$.
\emph{Case $j<h$, or $J'<r$ and $h\le j\le J'h$.} Put $e_j=s_j-q_j=-q_j\ind(i\le j)$, supported on at most $(r-1)h$ indices with entries at most $(2rh)^{r-1}/(r-1)!$.
\emph{Case $j>J'h$.} Put $e_j=s_j$, supported on $\{j+r,\dots,n\}$, at most $h-1$ indices, entries at most $(h+r-3)^{r-1}/(r-1)!$.

Since $\cP\subset V_J$, $(I-P_J)N=0$, and $\|(I-P_J)s_j\|_2\le\|e_j\|_2\le\sqrt{rh}\,c_\infty h^{r-1}$. For $E=(e_1,\dots,e_m)$, $(I-P_J)S=(I-P_J)E$; column sums of $|E|$ are at most $rh\,c_\infty h^{r-1}$ and row sums at most $3rh\,c_\infty h^{r-1}$, so the Schur test gives $\|E\|_{\op}\le\sqrt3\,r\,c_\infty h^r$.
\end{proof}

%----------------------------------------------------------------------
\subsection{Prior mass}

Let $A=n^{-1/2}\Pp S$, $B_\varrho=\{u\in\R^m:\|Au\|_2\le\varrho\}$ and $\varphi(\varrho)=-\log\Ql(B_\varrho)$.

\begin{lemma}[small ball]\label{lem:sb}
There is $C_{\rm sb}$ depending only on $r$ such that for all $n>r$, $\lambda>0$, $\varrho>0$,
$\varphi(\varrho)\le C_{\rm sb}\{1+(n^{r-1/2}/(\lambda\varrho))^{2/(2r-1)}\}$.
\end{lemma}

\begin{proof}
A Laplace$(\lambda)$ variable is $\sqrt Vg$ with $V$ exponential of mean $2/\lambda^2$ and $g\sim N(0,1)$ (Andrews and Mallows, 1974). Given $V=(V_j)$, $\|Au\|_2^2\sim\sum_q\mu_q\gamma_q^2$ with $\gamma_q$ i.i.d.\ $N(0,1)$ and $\mu_1\ge\mu_2\ge\dots$ the squared singular values of $G=A\,\mathrm{diag}(V_j^{1/2})$. On $\mathcal V=\{\sum_jV_j\le4m/\lambda^2\}$, which has probability at least $1/2$ by Markov's inequality:

\emph{Eigenvalues.} By Eckart--Young (the squared Frobenius error of the best approximation of $G$ of rank at most $J$ is $\sum_{q>J}\mu_q$) and \eqref{eq:F}, since $\Pp-(P_J-P)=I-P_J$ and $P_J-P$ has rank at most $J$,
$\sum_{q>J}\mu_q\le n^{-1}\sum_jV_j\|(I-P_J)s_j\|_2^2\le b_0J^{-(2r-1)}$ with $b_0=2^{2r+1}c_F^2n^{2r-1}/\lambda^2$; also $\sum_q\mu_q\le b_0$. Taking $J=\lfloor q/2\rfloor$ and using monotonicity, $\mu_q\le b\,q^{-2r}$ with $b=8\cdot6^{2r-1}c_F^2n^{2r-1}/\lambda^2$.

\emph{Gaussian small ball.} Let $K=\lceil(8b/\{(2r-1)\varrho^2\})^{1/(2r-1)}\rceil$. Then $\E\sum_{q>K}\mu_q\gamma_q^2\le\varrho^2/8$, so the tail is at most $\varrho^2/2$ with probability at least $3/4$. With $t_q^2=\varrho^2q^{2r}/(2bK)$, $\{|\gamma_q|\le t_q,\ q\le K\}$ forces the head to be at most $\varrho^2/2$, and $t_q\ge\vartheta_r(q/K)^r$ with $\vartheta_r=2/\sqrt{2r-1}$. Since $\pr(|\gamma|\le t)\ge0.48\min\{t,1\}$ and $K!\ge(K/e)^K$, the head has probability at least $e^{-\beta_rK}$, $\beta_r=r+\log(\max\{1,\sqrt{2r-1}/2\}/0.48)$. Hence $\Ql(B_\varrho)\ge\tfrac38e^{-\beta_rK}$. With $x=n^{r-1/2}/(\lambda\varrho)$, $K\le1+6(64c_F^2/(2r-1))^{1/(2r-1)}x^{2/(2r-1)}$.
\end{proof}

\begin{lemma}[quadratic shift]\label{lem:qshift}
Let $m\ge2$, $A$ any real matrix with columns $a_1,\dots,a_m$, $B_\varrho$ and $\varphi$ as above for this $A$, and $\delta_A=\max_{j<m}\|a_{j+1}-a_j\|_2$. If $\delta_A\le\lambda\varrho/2$, then for every $h\in\R^m$
\[
\Ql(h+B_\varrho)\ge\Ql(B_\varrho)\exp\{-c_K\lambda^2\|h\|_2^2\},\qquad c_K=\frac{1+K/2}{1-e^{-2}},\quad K=\frac{4\delta_A\varphi(\varrho/4)}{\lambda\varrho}.
\]
\end{lemma}

\begin{proof}
Let $B=B_\varrho$ and $\pi$ the density of $\Ql$. By symmetry, $\Ql(h+B)$ is the average of $\int_B\pi(w\pm h)dw$; the AM--GM inequality, $|w_j+h_j|+|w_j-h_j|=2\max\{|w_j|,|h_j|\}$ and Jensen give $\Ql(h+B)\ge\Ql(B)\exp\{-\lambda\sum_j\E_B(|h_j|-|w_j|)_+\}$, with $\E_B$ under $\Ql(\cdot\mid B)$. The marginal density $f_j$ of $w_j$ under $\Ql(\cdot\mid B)$ is symmetric and log-concave (Pr\'ekopa, 1973), so $\E_B(|h_j|-|w_j|)_+\le f_j(0)h_j^2$ and
$\Ql(h+B)\ge\Ql(B)\exp\{-\lambda\|h\|_2^2\max_jf_j(0)\}$.
It remains to show $f_j(0)\le c_K\lambda$. Let $j'$ be a neighbour of $j$, $A_{-j}$ the matrix without column $j$, $Q_{-j}$ the law of the other coordinates, $g(t)=Q_{-j}\{\|A_{-j}v+ta_j\|_2\le\varrho\}$ and $\Lambda(s)=\log Q_{-j}\{\|A_{-j}v\|_2\le s\}$. Shifting $v$ by $-te_{j'}$ gives $g(t)\ge e^{-\lambda|t|}e^{\Lambda(\varrho-|t|\delta_A)}$. $\Lambda$ is nondecreasing, nonpositive and concave (Pr\'ekopa, 1973; Borell, 1975), and $-\Lambda(\varrho/4)\le\varphi(\varrho/4)$ (sections of $B_{\varrho/4}$ have largest measure at $w_j=0$). Concavity gives $\Lambda(\varrho)-\Lambda(\varrho-a)\le4\varphi(\varrho/4)a/\varrho$ for $0\le a\le\varrho/2$. For $|t|\le1/\lambda$, $a=|t|\delta_A\le\varrho/2$, so $g(t)\ge g(0)e^{-\lambda|t|(1+K)}$ and
$\Ql(B)=\int(\lambda/2)e^{-\lambda|t|}g(t)dt\ge g(0)(1-e^{-2})/(2+K)$. Hence $f_j(0)=(\lambda/2)g(0)/\Ql(B)\le c_K\lambda$.
\end{proof}

For $w\ge r$ let $G_w=w^{-1}\sum_{k=0}^{w-1}\mathbf b(k/w)\mathbf b(k/w)^\top$, $\mathbf b(x)=(1,x,\dots,x^{r-1})^\top$, and $g_r=\inf_{w\ge r}\lambda_{\min}(G_w)\in(0,1]$ (a constant depending only on $r$), and $C_s=1+r/g_r$.

\begin{lemma}[smoothed truth]\label{lem:smooth}
For every integer $r\le w\le m$ there is $u_w\in\R^m$ with $\|u_w\|_2\le C_sV^{*}/(n^{r-1}\sqrt w)$ and $\nn{S(u^{*}-u_w)}\le C_sV^{*}(w/n)^{r-1/2}$.
\end{lemma}

\begin{proof}
For each $j$ take a window $W_j=\{a_j+1,\dots,a_j+w\}\ni j$, $x_l=(l-a_j-1)/w$ and the kernel $K^{(j)}(l)=w^{-1}\mathbf b(x_j)^\top G_w^{-1}\mathbf b(x_l)$ on $W_j$. It reproduces polynomials of degree at most $r-1$ at $j$ and $\|K^{(j)}\|_2\le r/(g_r\sqrt w)$. The vector $d_j=s_j-\sum_lK^{(j)}(l)s_l$ vanishes for $i\le a_j+1$ and, by reproduction of $p_i$, for $i\ge a_j+w+1$; so $\|d_j\|_2\le C_sw^{r-1/2}$. Let $u_w=\sum_ju^{*}_jK^{(j)}$ and use $\|u^{*}\|_1=V^{*}n^{-(r-1)}$.
\end{proof}

\begin{proposition}[prior mass]\label{prop:prior}
Let $V^{*}\le\Vb$. There is $n_1$ depending only on $(c,\Vb,r)$ such that for $n\ge n_1$ there is $u^\circ$ with
$\nn{\Pp S(u^{*}-u^\circ)}\le\en/2$ and $-\log\Ql(u^\circ+B_{\en/2})\le C_{\rm pr}\,n\en^2$, where
$C_{\rm pr}=C_{\rm sb}\{1+(2/c)^{2/(2r-1)}\}+4c^2C_s^2\Vb^2(2C_s\Vb)^{2/(2r-1)}$.
\end{proposition}

\begin{proof}
Note $\en^{-2/(2r-1)}=n\en^2$. If $\Vb=0$ take $u^\circ=0$ and use Lemma~\ref{lem:sb}. Otherwise let $w=\lfloor n(\en/(2C_s\Vb))^{2/(2r-1)}\rfloor$ and $u^\circ=u_w$ (Lemma~\ref{lem:smooth}); for $n\ge n_1$, $r\le w\le m$ and the approximation error is at most $\en/2$, while
$\lambda^2\|u^\circ\|_2^2\le c^2C_s^2\Vb^2n/w\le2c^2C_s^2\Vb^2(2C_s\Vb)^{2/(2r-1)}n\en^2$.
Apply Lemma~\ref{lem:qshift} with $A=n^{-1/2}\Pp S$, $\varrho=\en/2$: by \eqref{eq:pascal}, $\delta_A\le n^{r-2}$, so $\delta_A\le\lambda\varrho/2$ and, by Lemma~\ref{lem:sb}, $K\to0$ and $c_K\le2$ for $n\ge n_1$. Hence $-\log\Ql(u^\circ+B_{\en/2})\le\varphi(\en/2)+2\lambda^2\|u^\circ\|_2^2\le C_{\rm pr}n\en^2$.
\end{proof}

%----------------------------------------------------------------------
\subsection{Denominator}

\begin{lemma}[denominator]\label{lem:evid}
Let $n\ge n_1$ and $C_E=\tfrac54C_1+1+C_{\rm pr}$. There is an event $\Omega_E$ with $\pr(\Omega_E^c)\le\tfrac54(n\en^2)^{-1}$ on which $N(\R^n)\ge V_r(\sqrt n\,\en/2)^re^{-C_En\en^2}$.
\end{lemma}

\begin{proof}
We adapt Ghosal, Ghosh and van der Vaart (2000, Lemma~8.1). Let $\mathcal G=\{Nc+Su:\ u\in u^\circ+B_{\en/2},\ \|c-\bar c(u)\|_2\le\sqrt n\en/2\}$. By \eqref{eq:pyth}, $\nn{\delta}^2\le\tfrac54\en^2$ on $\mathcal G$, and by \eqref{eq:vol}, $\nu(\mathcal G)=V_r(\sqrt n\en/2)^r\Ql(u^\circ+B_{\en/2})$. Let $\bar\nu=\nu(\cdot\cap\mathcal G)/\nu(\mathcal G)$ and $\bar{\mathcal L}=\int\Lh\,d\bar\nu$. Jensen gives $N(\R^n)\ge\nu(\mathcal G)e^{-\bar{\mathcal L}}$. By Lemma~\ref{lem:check}(d), $\E\bar{\mathcal L}\le\tfrac54C_1n\en^2$; by Lemma~\ref{lem:check}(a) and independence, $\mathrm{Var}(\bar{\mathcal L})\le\int\|\delta\|_2^2d\bar\nu\le\tfrac54n\en^2$. Chebyshev's inequality gives $\pr\{\bar{\mathcal L}\ge(\tfrac54C_1+1)n\en^2\}\le\tfrac54(n\en^2)^{-1}$. Use Proposition~\ref{prop:prior}.
\end{proof}

%----------------------------------------------------------------------
\subsection{Localization}

For $\rho,T>0$ let $K(\rho,T)=\{\delta:\ \nn{\delta}\le\rho,\ \TV^{(r)}(\delta)\le T\}$, a compact convex symmetric set.

\begin{lemma}[width]\label{lem:width}
For independent Rademacher signs $\xi_i$, all $\rho,T>0$ and $1\le J\le m$,
$\E\sup_{K(\rho,T)}\langle\xi,\delta\rangle\le\sqrt n\{\rho\sqrt{J+r}+2^rc_F\sqrt{\log(2n)}\,T\,J^{-(r-1/2)}\}$.
\end{lemma}

\begin{proof}
For $\delta=Nc+Su$, $\|u\|_1\le Tn^{-(r-1)}$ and $\langle\xi,\delta\rangle=\langle P_J\xi,\delta\rangle+\langle S^\top(I-P_J)\xi,u\rangle$. Bound the first term by $\|P_J\xi\|_2\sqrt n\rho$ with $\E\|P_J\xi\|_2\le\sqrt{J+r}$, and the second by $\max_j|\langle(I-P_J)s_j,\xi\rangle|\,Tn^{-(r-1)}$, a maximum of $2m$ sub-Gaussian variables with variance proxy $c_F^2h^{2r-1}$ by \eqref{eq:F}. Use $h\le2n/J$.
\end{proof}

\begin{lemma}[interpolation]\label{lem:interp}
There is $C_I$ depending only on $r$ such that for $n>r$ and $\delta\in\R^n$,
$\|\delta\|_\infty\le C_I\{\nn{\delta}+\nn{\delta}^{(2r-2)/(2r-1)}\TV^{(r)}(\delta)^{1/(2r-1)}\}$.
\end{lemma}

\begin{proof}
Let $a=\nn{\delta}$, $T=\TV^{(r)}(\delta)$, $s=|\delta_{i_0}|=\|\delta\|_\infty$, $C_0=r/g_r$. On a window $W\ni i_0$ of $w\ge r$ consecutive indices, Lemma~\ref{lem:coord} applied to $\delta_W$ gives $\delta_W=p+S_wu_W$ with $p$ a polynomial and $\|\delta_W-p\|_\infty\le\varrho_w=w^{r-1}T/\{(r-1)!n^{r-1}\}$. Since $\max_W|p|\le(C_0/w)^{1/2}\|p\|_{\ell_2(W)}$ and $\|p\|_{\ell_2(W)}\le\sqrt na+\sqrt w\varrho_w$,
$s\le\sqrt{C_0}\,a\sqrt{n/w}+C'T(w/n)^{r-1}$ with $C'=(1+\sqrt{C_0})/(r-1)!$. If $T\le a$ take $w=n$. Otherwise let $x^{*}=(a/T)^{2/(2r-1)}$; if $nx^{*}\ge r$ take $w=\lceil nx^{*}\rceil$, and if $nx^{*}<r$ use $s\le\sqrt na$. The claim holds with $C_I=\sqrt{C_0}+2^{r-1}C'$.
\end{proof}

\begin{lemma}[concentration]\label{lem:conc}
Let $K\subset\R^n$ be compact and symmetric and $G_K=\sup_{\delta\in K}|\Lh(\delta)-\Lm(\delta)|$. Then
$\E G_K\le4\E\sup_K\langle\xi,\delta\rangle$ and $\mathrm{Var}(G_K)\le\sup_K\|\delta\|_2^2$.
\end{lemma}

\begin{proof}
Symmetrization for independent processes, in absolute-value form (van der Vaart and Wellner, 1996, Lemma~2.3.6), gives $\E G_K\le2\E\sup_K|\sum_i\xi_i\ell_i(\delta_i)|$; the $\ell_i$ are contractions vanishing at $0$, so the contraction principle in absolute-value form (Ledoux and Talagrand, 1991, Theorem~4.12) and the symmetry of $K$ give the factor $4$. For the variance, replacing $Z_i$ by an independent copy changes $G_K$ by at most $|\hat\delta_i|$, with $\hat\delta$ a maximizer (Lemma~\ref{lem:check}(a)); the Efron--Stein inequality (Efron and Stein, 1981; Boucheron, Lugosi and Massart, 2013, Theorem~3.1) gives $\mathrm{Var}(G_K)\le\|\hat\delta\|_2^2$.
\end{proof}

Let
\[
\rho_n^2=n^{-(2r-1)/(2r+1)}\log n,\qquad T_n=\frac{4n^r}{\lambda}=\frac4c\sqrt n,\qquad K_n=K(\rho_n,T_n),\qquad \alpha_n=\frac{\fl}4n\rho_n^2,
\]
and $s_0=\min\{1,L,3\fl/(4\fu)\}$.

\begin{proposition}[localization]\label{prop:local}
There is $n_2\ge n_1$, depending only on $(c,r,L,\fl,\fu,\Vb)$, and for $n\ge n_2$ an event $\Omega_G$ with $\pr(\Omega_G^c)\le64/(\fl^2n\rho_n^2)$, such that
\begin{enumerate}[label=(\alph*),itemsep=1pt,topsep=2pt]
\item $\sup_{\delta\in K_n}\|\delta\|_\infty\le s_0$; in particular $\Delta_n^2(\delta)=\nn{\delta}^2$ on $K_n$;
\item on $\Omega_G\cap\Omega_E$, $\Pi_\lambda(\theta-\thetas\notin K_n\mid y)\le r!\{e^{-\fl n\rho_n^2/16}+e^{-n/8}\}$.
\end{enumerate}
\end{proposition}

\begin{proof}
(a) By Lemma~\ref{lem:interp}, $\sup_{K_n}\|\delta\|_\infty\le C_I\{\rho_n+(\rho_n^{2r-2}T_n)^{1/(2r-1)}\}$, and
\[
\rho_n^{2r-2}T_n=\frac4c\,n^{-(4r^2-8r+1)/(2(2r+1))}(\log n)^{r-1}\to0
\]
because $4r^2-8r+1\ge1$ for $r\ge2$. This is the substantive use of $r\ge2$: at $r=1$ the exponent is $+1/2$. (Case~I of Lemma~\ref{lem:F} and the bound $\delta_A\le n^{r-2}$ in Proposition~\ref{prop:prior} also use $r\ge2$, but have elementary replacements at $r=1$.)

\emph{The event.} Let $G=G_{K_n}$ and $J_n=\lceil(T_n\sqrt{\log(2n)}/\rho_n)^{1/r}\rceil$. Lemma~\ref{lem:width} gives $\E\sup_{K_n}\langle\xi,\delta\rangle\le(1+2^rc_F)\sqrt n\rho_n\sqrt{J_n+r}$, and $J_n+r=O(n^{2/(2r+1)})=o(n\rho_n^2)$, so $\E G\le(\fl/8)n\rho_n^2$ for $n\ge n_2$. Let $\Omega_G=\{G\le(\fl/4)n\rho_n^2\}$; Lemma~\ref{lem:conc} and Chebyshev give the bound on $\pr(\Omega_G^c)$.

\emph{Rays.} Let $t(\delta)=\max\{\nn{\delta}/\rho_n,\TV^{(r)}(\delta)/T_n\}$ and, for $\delta\notin K_n$, $\hat\delta=\delta/t(\delta)$; by convexity $\Lh(\delta)\ge t(\delta)\Lh(\hat\delta)$. On $\Omega_G$: if $\nn{\delta}/\rho_n\ge\TV^{(r)}(\delta)/T_n$ (set $A_1$), then $\|\hat\delta\|_\infty\le s_0$, Lemma~\ref{lem:check}(d) gives $\Lm(\hat\delta)\ge(\fl/2)n\rho_n^2$, and $\Lh(\delta)\ge\alpha_n\nn{\delta}/\rho_n$; otherwise (set $A_2$), $\Lh(\delta)\ge-\alpha_nt(\delta)$.

\emph{Mass of $A_1$.} On $A_1$, $\nn{\delta}\ge\rho_n/2+\|c-\bar c(u)\|_2/(2\sqrt n)$; with \eqref{eq:vol} and $\int_{\R^r}e^{-\kappa\|c\|_2}dc=r!V_r\kappa^{-r}$, $N(A_1)\le r!V_r(2\sqrt n\rho_n/\alpha_n)^re^{-\alpha_n/2}$.

\emph{Mass of $A_2$.} $\Gamma=\lambda\|u\|_1\sim\mathrm{Gamma}(m,1)$ under $\Ql$, and $t(\delta)\le\bar t:=\Gamma/(4n)+c\Vb/(4\sqrt n)\le\Gamma/(4n)+1/2$, so $A_2\subset\{\Gamma>2n\}$. By \eqref{eq:vol}, $x^r\le r!e^x$, $\ind\{\Gamma>2n\}\le e^{(\Gamma-2n)/2}$ and $\E e^{\vartheta\Gamma}=(1-\vartheta)^{-m}$,
$N(A_2)\le r!V_r(\sqrt n\rho_n)^r\E_{\Ql}[e^{(\alpha_n+1)\bar t}\ind\{\Gamma>2n\}]\le r!V_r(\sqrt n\rho_n)^re^{-n/4}$ for $n\ge n_2$.

(b) Divide by Lemma~\ref{lem:evid}: $\Pi_\lambda(K_n^c\mid y)\le r!(16/(\fl n\rho_n\en))^re^{-\alpha_n/2+C_En\en^2}+r!(2\rho_n/\en)^re^{-n/4+C_En\en^2}$. Since $n\en^2=o(n\rho_n^2)$ and $\rho_n/\en$ grows polynomially, this is at most $r!\{e^{-\alpha_n/4}+e^{-n/8}\}$ for $n\ge n_2$.
\end{proof}

%----------------------------------------------------------------------
\subsection{Shells}

For $1\le J\le m$ let $t=S^\top(I-P_J)\psi$, i.e.\ $t_j=\langle(I-P_J)s_j,\psi\rangle$, and
\[
\Omega_\psi(J)=\Bigl\{\|P_J\psi\|_2^2\le J+r,\ \ \|t\|_\infty\le c_Fh^{r-1/2}\sqrt{\log(2n)},\ \ \|t\|_2^2\le c_F^2mh^{2r-1}\Bigr\}.
\]

\begin{lemma}[score events]\label{lem:score}
For every $1\le J\le m$, $\pr(\Omega_\psi(J)^c)\le\frac4{9(J+r)}+\frac1{2n}+\frac{8r}{3J}$.
\end{lemma}

\begin{proof}
The $\psi_i$ are i.i.d., centred, $|\psi_i|\le1$, variance $v\le1/4$, so for symmetric $Q$: $\E\psi^\top Q\psi=v\,\mathrm{tr}Q$ and $\mathrm{Var}(\psi^\top Q\psi)\le\frac14\|Q\|_F^2$. With $Q=P_J$, Chebyshev gives the first term. Hoeffding's inequality and \eqref{eq:F} give $\pr(|t_j|>c_Fh^{r-1/2}\sqrt{\log(2n)})\le1/(2n^2)$; sum over $j$. For $W=(I-P_J)S$, $\|W\|_F^2\le b:=c_F^2mh^{2r-1}$ and $\|WW^\top\|_F\le c_{\op}h^r\sqrt b$ by \eqref{eq:F}, so $\pr(\|t\|_2^2>b)\le(c_{\op}^2h^{2r}b/4)/(3b/4)^2=\frac{4r}3\frac hm\le\frac{8r}{3J}$.
\end{proof}

For $a>0$ let $E_a=\{\theta:\ \|\delta\|_\infty\le s_0,\ an\en^2<\|\delta\|_2^2\le2an\en^2\}$.

\begin{proposition}[shells]\label{prop:shell}
Let $M>0$, $J=\lfloor(\fl^2/128)Mn\en^2\rfloor-r$ and $C_{\rm fine}=2^{2r+1}c_F^2(256/\fl^2)^{2r-1}$. There is $n_3$, depending only on $(M,c,r,\fl,\Vb)$, such that for $n\ge n_3$: $1\le J\le m$, $J\ge(\fl^2/256)Mn\en^2$, and for every $a\ge M$
\[
\E\bigl[N(E_a)\ind_{\Omega_\psi(J)}\bigr]\le V_r(2an\en^2)^{r/2}\exp\Bigl\{-\frac{\fl}8an\en^2+\Bigl(1+\frac{C_{\rm fine}}{c^2M^{2r-1}}\Bigr)n\en^2\Bigr\}.
\]
\end{proposition}

\begin{proof}
By Lemma~\ref{lem:check}(c) and Cauchy--Schwarz in $L^2(\nu)$,
$N(E_a)\le I_1^{1/2}I_2^{1/2}$ with $I_1=\int_{E_a}e^{2\langle\psi,\delta\rangle}d\nu$ and $I_2=\int_{E_a}e^{-2\Rh(\delta)}d\nu$, $\Rh=\sum_iR_i(\delta_i)$.

\emph{Remainder.} On $E_a$, $\|\delta\|_\infty\le s_0$, so Lemma~\ref{lem:check}(e), (d) and independence give $\E e^{-2\Rh(\delta)}\le e^{-\Lm(\delta)}\le e^{-(\fl/2)an\en^2}$; by Tonelli and \eqref{eq:vol}, $\E I_2\le V_r(2an\en^2)^{r/2}e^{-(\fl/2)an\en^2}$.

\emph{Score, on $\Omega_\psi=\Omega_\psi(J)$.} Since $(I-P_J)N=0$, $\langle\psi,\delta\rangle=\langle P_J\psi,\delta\rangle+\langle t,u-u^{*}\rangle$. On $E_a$ with $a\ge M$, $|\langle P_J\psi,\delta\rangle|\le\sqrt{J+r}(2an\en^2)^{1/2}\le(\fl/8)an\en^2$. By \eqref{eq:vol} and $\int e^{sx}(\lambda/2)e^{-\lambda|x|}dx=(1-s^2/\lambda^2)^{-1}$,
\[
I_1\le e^{(\fl/4)an\en^2}e^{-2\langle t,u^{*}\rangle}V_r(2an\en^2)^{r/2}\prod_j(1-4t_j^2/\lambda^2)^{-1}.
\]
Using $h\le2n/J$, $J\ge(\fl^2/256)Mn\en^2$, $\lambda^2=c^2n^{2r-1}$, $(n\en^2)^{2r}=n$ and $\|u^{*}\|_1\le\Vb n^{-(r-1)}$: for $n\ge n_3$, $4\|t\|_\infty^2/\lambda^2\le2^{2r+1}c_F^2\log(2n)/(c^2J^{2r-1})\le\frac12$; $8\|t\|_2^2/\lambda^2\le2C_{\rm fine}c^{-2}M^{-(2r-1)}n\en^2$; and $2|\langle t,u^{*}\rangle|\le2c_F\Vb(512/(\fl^2M))^{r-1/2}\sqrt{\log(2n)}(n\en^2)^{1/2}\le2n\en^2$. With $-\log(1-x)\le2x$ on $[0,\frac12]$,
$I_1\le\bar I_1:=V_r(2an\en^2)^{r/2}\exp\{(\fl/4)an\en^2+2n\en^2+2C_{\rm fine}c^{-2}M^{-(2r-1)}n\en^2\}$ on $\Omega_\psi$.

Since $\bar I_1$ is deterministic, $\E[N(E_a)\ind_{\Omega_\psi}]\le\bar I_1^{1/2}(\E I_2)^{1/2}$ by Jensen.
\end{proof}

The score is never averaged over the data: only the nonnegative remainder is. This is why no lower bound on $\sigma_0$ is needed.

%----------------------------------------------------------------------
\subsection{Proof of Theorem~\ref{thm:main}(a)}

Let
\[
M=\frac{32}{\fl}\bigl(C_E+1\bigr)+\Bigl(\frac{32C_{\rm fine}}{\fl c^2}\Bigr)^{1/(2r)},\qquad C=2+\frac{768r}{\fl^2M},
\]
$J$ as in Proposition~\ref{prop:shell}, $n\ge n_0:=\max\{n_2,n_3\}$ (enlarged finitely often below) and $a_k=2^kM$. If $\delta\in K_n$ and $\Delta_n^2(\delta)>M\en^2$, then by Proposition~\ref{prop:local}(a) $\nn{\delta}^2>M\en^2$ and $\theta\in E_{a_k}$ for some $k\ge0$. Hence, with $\Omega_\psi=\Omega_\psi(J)$,
\begin{align*}
\E\,\Pi_\lambda(\Delta_n^2>M\en^2\mid y)\le{}&\pr(\Omega_E^c)+\pr(\Omega_G^c)+\pr(\Omega_\psi^c)\\
&+\E[\Pi_\lambda(K_n^c\mid y)\ind_{\Omega_E\cap\Omega_G}]+\sum_{k\ge0}\E[\Pi_\lambda(E_{a_k}\mid y)\ind_{\Omega_E\cap\Omega_\psi}].
\end{align*}
\emph{Shells.} On $\Omega_E$ the denominator is at least the deterministic bound of Lemma~\ref{lem:evid}; the volume factors leave $(2an\en^2)^{r/2}(\sqrt n\en/2)^{-r}=(8a)^{r/2}$, and Proposition~\ref{prop:shell} gives
\[
\E[\Pi_\lambda(E_a\mid y)\ind_{\Omega_E\cap\Omega_\psi}]\le(8a)^{r/2}\exp\Bigl\{-\frac{\fl}8an\en^2+\Bigl(C_E+1+\frac{C_{\rm fine}}{c^2M^{2r-1}}\Bigr)n\en^2\Bigr\}.
\]
By the choice of $M$, $(\fl/32)M\ge C_E+1$ and $(\fl/32)M^{2r}\ge C_{\rm fine}/c^2$, so the bracket is at most $(\fl/16)M\le(\fl/16)a$. With $x_n=(\fl/16)Mn\en^2\ge r$ and $2^{kr/2}e^{-(2^k-1)x_n}\le e^{-0.65kr}$, the shell terms sum to at most $2(8M)^{r/2}e^{-x_n}$.

\emph{Conclusion.} By Lemma~\ref{lem:evid}, Proposition~\ref{prop:local}, Lemma~\ref{lem:score} with $J\ge(\fl^2/256)Mn\en^2$ and $4/9+8r/3\le3r$, the right side is at most
\[
\frac5{4n\en^2}+\frac{64}{\fl^2n\rho_n^2}+\frac{768r}{\fl^2Mn\en^2}+\frac1{2n}+r!\{e^{-\fl n\rho_n^2/16}+e^{-n/8}\}+2(8M)^{r/2}e^{-x_n}\le C\,n^{-1/(2r)},
\]
since $n\en^2=n^{1/(2r)}$ and every term other than the first and third is $o(1/(n\en^2))$. \qed

%----------------------------------------------------------------------
\subsection{Proof of Theorem~\ref{thm:main}(b)}

\begin{lemma}[symmetric log-concave densities]\label{lem:lc}
A symmetric log-concave density $p$ on $\R$ with variance $v\in(0,\infty)$ has $\sup p=p(0)\le1/\sqrt{2v}$.
\end{lemma}

\begin{proof}
Let $q(x)=(\beta/2)e^{-\beta|x|}$ with $\beta=2p(0)$. On $[0,\infty)$, $\log p-\log q$ is concave and vanishes at $0$, so $p\ge q$ exactly on some $[0,x_0]$. As $\int_0^\infty(p-q)=0$, $\int_0^\infty x^2(p-q)=\int_0^\infty(x^2-x_0^2)(p-q)\le0$, so $v\le2/\beta^2$.
\end{proof}

\begin{lemma}[small ball, upper bound]\label{lem:sbu}
Let $C_\ell=8e^2(r+1)\binom{2r}r^{1/2}$, $x_0=C_\ell(4(r+1))^{r-1/2}$, $c_{\rm lb}=C_\ell^{-2/(2r-1)}/(4(r+1))$. If $x=n^{r-1/2}/(\lambda\varrho)$ satisfies $x_0\le x\le C_\ell n^{r-1/2}$, then $\Ql(B_\varrho)\le\exp\{-c_{\rm lb}x^{2/(2r-1)}\}$.
\end{lemma}

\begin{proof}
Let $\ell\ge1$ with $2(r+1)\ell\le n$ and $\mathsf B=\lfloor n/((r+1)\ell)\rfloor$. Let $\kappa$ be the coefficient vector of $(1+z+\dots+z^{\ell-1})^{r+1}$, so $\|\kappa\|_2^2\ge\ell^{2r+1}/(r+1)$; let $\kappa_b$ be $\kappa$ placed at offset $j_b=(b-1)(r+1)\ell$, and $g_b=D^\top\kappa_b$. The generating function of $g_b$ is $z^{j_b+1}(z^\ell-1)^r\sigma_\ell(z)$, $\sigma_\ell(z)=1+\dots+z^{\ell-1}$, so the $g_b$ have disjoint supports and $\|g_b\|_2^2=\ell\binom{2r}r$. By Lemma~\ref{lem:coord}(a), $\langle g_b,\Pp Su\rangle=\langle\kappa_b,u\rangle=:Y_b$. For $u\in B_\varrho$, $\sum_bY_b^2/\|g_b\|_2^2\le n\varrho^2$, so at least $\lceil\mathsf B/2\rceil$ of the $Y_b$ satisfy $Y_b^2\le y_0^2=2n\varrho^2\ell\binom{2r}r/\mathsf B$. Under $\Ql$ the $Y_b$ are independent, symmetric and log-concave with variance $2\|\kappa\|_2^2/\lambda^2$, so by Lemma~\ref{lem:lc}
$p_0=\Ql(|Y_b|\le y_0)\le2(r+1)\binom{2r}r^{1/2}x^{-1}(n/\ell)^{r-1/2}$,
and a union bound gives $\Ql(B_\varrho)\le(4p_0)^{\mathsf B/2}$. Take $\ell=\lceil n(C_\ell/x)^{2/(2r-1)}\rceil$: then $4p_0\le e^{-2}$, the range of $x$ gives $2(r+1)\ell\le n$, and $\mathsf B\ge(x/C_\ell)^{2/(2r-1)}/(4(r+1))$.
\end{proof}

\begin{proof}[Proof of (b)]
Let $\eta=\sqrt{c_0}\,\en$, $c_0\le1$, and $H=\{\nn{\delta}\le\eta\}$. By Proposition~\ref{prop:local}(a), $\{\Delta_n^2(\delta)\le c_0\en^2\}\subset K_n^c\cup H$. By Lemma~\ref{lem:check}(f), Tonelli, \eqref{eq:pyth}, \eqref{eq:vol} and Anderson's inequality (Anderson, 1955; $\Ql$ has a symmetric log-concave density),
\[
\E N(H)\le e^{n\eta^2/8}\nu(H)\le e^{n\eta^2/8}V_r(\sqrt n\eta)^r\Ql(u^{*}+B_\eta)\le e^{n\eta^2/8}V_r(\sqrt n\eta)^r\Ql(B_\eta).
\]
Here $x=1/(c\sqrt{c_0}\en)\in[x_0,C_\ell n^{r-1/2}]$ for large $n$ and $x^{2/(2r-1)}=(c^2c_0)^{-1/(2r-1)}n\en^2$, so Lemma~\ref{lem:sbu} gives $\Ql(B_\eta)\le\exp\{-c_{\rm lb}(c^2c_0)^{-1/(2r-1)}n\en^2\}$. With Lemma~\ref{lem:evid},
\[
\E[\Pi_\lambda(H\mid y)\ind_{\Omega_E}]\le2^r\exp\bigl\{\bigl(\tfrac18+C_E-c_{\rm lb}(c^2c_0)^{-1/(2r-1)}\bigr)n\en^2\bigr\}\le2^re^{-n\en^2}
\]
once $c_0$ is so small that $c_{\rm lb}(c^2c_0)^{-1/(2r-1)}\ge C_E+\frac98$. Adding $\pr(\Omega_E^c)+\pr(\Omega_G^c)+\E[\Pi_\lambda(K_n^c\mid y)\ind_{\Omega_E\cap\Omega_G}]$ (Lemma~\ref{lem:evid}, Proposition~\ref{prop:local}), the total is at most $2/(n\en^2)=2n^{-1/(2r)}$ for large $n$.
\end{proof}

%======================================================================
\begin{thebibliography}{99}\small\setlength{\itemsep}{0pt}
\bibitem{anderson} Anderson, T. W. (1955). The integral of a symmetric unimodal function over a symmetric convex set and some probability inequalities. \emph{Proc. Amer. Math. Soc.} 6, 170--176.
\bibitem{am} Andrews, D. F. and Mallows, C. L. (1974). Scale mixtures of normal distributions. \emph{J. Roy. Statist. Soc. B} 36, 99--102.
\bibitem{atchade} Atchad\'e, Y. A. (2017). On the contraction properties of some high-dimensional quasi-posterior distributions. \emph{Ann. Statist.} 45, 2248--2273.
\bibitem{bc} Belloni, A. and Chernozhukov, V. (2011). $\ell_1$-penalized quantile regression in high-dimensional sparse models. \emph{Ann. Statist.} 39, 82--130.
\bibitem{borell} Borell, C. (1975). Convex set functions in $d$-space. \emph{Period. Math. Hungar.} 6, 111--136.
\bibitem{blm} Boucheron, S., Lugosi, G. and Massart, P. (2013). \emph{Concentration Inequalities}. Oxford University Press.
\bibitem{dj} Donoho, D. L. and Johnstone, I. M. (1998). Minimax estimation via wavelet shrinkage. \emph{Ann. Statist.} 26, 879--921.
\bibitem{es} Efron, B. and Stein, C. (1981). The jackknife estimate of variance. \emph{Ann. Statist.} 9, 586--596.
\bibitem{ggv} Ghosal, S., Ghosh, J. K. and van der Vaart, A. W. (2000). Convergence rates of posterior distributions. \emph{Ann. Statist.} 28, 500--531.
\bibitem{knight} Knight, K. (1998). Limiting distributions for $L_1$ regression estimators under general conditions. \emph{Ann. Statist.} 26, 755--770.
\bibitem{lt} Ledoux, M. and Talagrand, M. (1991). \emph{Probability in Banach Spaces}. Springer.
\bibitem{padilla} Madrid Padilla, O. H. and Chatterjee, S. (2022). Risk bounds for quantile trend filtering. \emph{Biometrika} 109, 751--768.
\bibitem{prekopa} Pr\'ekopa, A. (1973). On logarithmic concave measures and functions. \emph{Acta Sci. Math. (Szeged)} 34, 335--343.
\bibitem{srg} Sriram, K., Ramamoorthi, R. V. and Ghosh, P. (2013). Posterior consistency of Bayesian quantile regression based on the misspecified asymmetric Laplace density. \emph{Bayesian Anal.} 8, 479--504.
\bibitem{sm} Syring, N. and Martin, R. (2023). Gibbs posterior concentration rates under sub-exponential type losses. \emph{Bernoulli} 29, 1080--1108.
\bibitem{tibs} Tibshirani, R. J. (2014). Adaptive piecewise polynomial estimation via trend filtering. \emph{Ann. Statist.} 42, 285--323.
\bibitem{vw} van der Vaart, A. W. and Wellner, J. A. (1996). \emph{Weak Convergence and Empirical Processes}. Springer.
\bibitem{ywh} Yang, Y., Wang, H. J. and He, X. (2016). Posterior inference in Bayesian quantile regression with asymmetric Laplace likelihood. \emph{Int. Statist. Rev.} 84, 327--344.
\end{thebibliography}

\end{document}
